\begin{textsample}{POS Dim 3 – 2009\_c6 – Score -2.00 – 2009\_c6/enade\_2009\_secr\_exec\_8db19b640f.txt}  \label{ex:f3_pos_240}
Capítulo 3 Análise Técnica da Prova Este capítulo tem por objetivo apresentar o desempenho dos estudantes de SECRETARIADO EXECUTIVO no ENADE/2009.
Para isso, foram calculadas as estatísticas básicas da prova como um todo, bem como as estatísticas das questões relacionadas à Formação Geral e ao Componente Específico.
Nas Tabelas, são apresentadas as seguintes estatísticas : total da população, de presentes, média, erro-padrão da média, desvio-padrão, mediana, nota máxima, nota mínima e coeficiente de assimetria ( ver definição em 1.4.3 ).
Nas estatísticas das questões discursivas, foram detalhadas também as notas nulas : em branco, por protesto ou por insuficiência.
As estatísticas apresentadas neste capítulo contemplaram, separadamente, os ingressantes, os concluintes e o total de estudantes, e foram calculadas tendo -se em vista as seguintes agregações : ( a ) Região Geográfica e País e ( b ) a Categoria Administrativa e a Organização Acadêmica.
As notas brutas das questões discursivas em cada uma das partes, Formação Geral/Componente Específico, foram calculadas a partir da média simples das notas de cada uma das questões correspondentes.
A primeira questão discursiva de Formação Geral foi desconsiderada e, portanto, a nota bruta referente às questões discursivas de Formação Geral foi igual à nota da questão mantida, a segunda.
No Componente Específico, a nota bruta correspondeu à média das notas nas três questões.
As notas brutas das questões objetivas de cada uma das partes, Formação Geral/Componente Específico, foram calculadas a partir dos acertos dos estudantes considerando -se uma regra de três para a conversão para a escala de 0 a 100, descontando -se aquelas questões que foram desconsideradas pela comissão ou que não passaram pelo teste do índice de discriminação ponto-bisserial ( ver seção 1.4.6 para a descrição do índice ).
As questões com índices de discriminação inferior a 0,2 foram excluídas do cômputo da nota do estudante.
As notas brutas de Formação Geral e do Componente Específico para cada um dos estudantes foram calculadas como médias ponderadas das suas notas nas questões objetivas e discursivas, ponderação diferente para cada parte : C\_FG = 0,6 × C\_FG\textasciicircum{}OBJ + 0,4 × C\_FG\textasciicircum{}DIS C\_CE = 0,85 × C\_CE\textasciicircum{}OBJ + 0,15 × C\_CE\textasciicircum{}DIS, onde C\_FG é a nota bruta de Formação Geral, C\_FG\textasciicircum{}OBJ é a nota bruta nas questões objetivas de Formação Geral, C\_FG\textasciicircum{}DIS é a nota bruta nas questões discursivas de Formação Geral, C\_CE é a nota bruta do Componente Específico, C\_CE\textasciicircum{}OBJ é a nota bruta nas questões objetivas do Componente Específico, e C\_CE\textasciicircum{}DIS é a nota bruta nas questões discursivas do Componente Específico.
A nota bruta da prova, C\_GER, por sua vez, é também uma média ponderada da nota de cada uma das partes, Formação Geral/Componente Específico : C\_GER = 0,25 × C\_FG + 0,75 × C\_CE. 3.1 Estatísticas Básicas da Prova 3.1.1 Estatísticas Básicas Gerais A Tabela 3.1 apresenta as estatísticas básicas da prova por grupo de estudantes.
O total de estudantes inscritos para a prova foi de 5620.
Destes, 77,5 \% compareceram, sendo que a presença foi maior entre concluintes ( 88,0 \% ) do que entre ingressantes ( 68,6 \% ).
A média geral da prova foi de 41,19, sendo que os ingressantes obtiveram média mais baixa, de 39,27, que os concluintes, cuja média foi 42,95.
O desvio-padrão geral foi de 14,4, sendo o do grupo de ingressantes menor ( 13,8 ), do que o do grupo dos concluintes ( 14,7 ), aqueles possuíram uma distribuição mais homogênea nas notas.
A nota máxima foi 90,2, obtida por pelo menos um concluinte ao passo que a maior nota obtida por um ingressante foi 85,7.
Os coeficientes de assimetria para os dois grupos considerados, o de concluintes e o de ingressantes, bem como para a população como um todo são próximos de zero, indicando uma maior simetria.
Uma comparação entre o desempenho de ingressantes e concluintes pode ser realizada a partir dos dados apresentados no Gráfico 3.1.
Cumpre destacar que as notas dos concluintes, de um modo geral, foram mais elevadas do que as dos ingressantes, cuja distribuição de notas está à esquerda dos primeiros.
A ideia de melhor desempenho é reforçada pelas notas medianas dos dois grupos : 42,6 entre os concluintes e 38,8 entre os ingressantes.
Nas faixas de notas até 40 pontos o percentual de ingressantes é maior que o dos concluintes, situação esta que se inverte para as faixas de notas superiores a 50 pontos.
Constata -se que aproximadamente 22,6 \% dos ingressantes e 31,6 \% dos concluintes situaram -se nas faixas de notas superiores a 50,0 pontos.
O Gráfico 3.2 apresenta informações referentes ao desempenho geral de ingressantes e concluintes, comparando os resultados considerando Região, Categoria Administrativa e Organização Acadêmica.
Levando -se em conta as notas médias dos estudantes em cada Região, observa -se que há, nos dois grupos, uma variação expressiva entre as Regiões.
As notas médias mais elevadas entre os concluintes ocorreram nas Regiões Nordeste ( 43,03 ), Sudeste ( 45,74 ) e Sul ( 44,42 ), ambas acima da média nacional ( 42,95 ).
Em relação aos ingressantes, as médias mais altas foram obtidas nas Regiões Sudeste ( 41,90 ), Norte ( 40,80 ) e Sul ( 39,52 ) nas demais Regiões a média foi inferior à média nacional ( 39,27 ).
Na Região Centro-Oeste foi obtida a média de notas mais baixa entre os ingressantes ( 31,90 ) assim como entre os concluintes ( 35,84 ).
O Norte apresentou a média dos ingressantes superior a dos concluintes.
Cursos nesta área foram oferecidos em instituições federais, estaduais, municipais e privadas.
Levando -se em conta os agrupamentos dos estudantes em categorias administrativas, observa -se que a pontuação média mais elevada entre os concluintes e ingressantes foi encontrada em instituições de origem estadual ( respectivamente 48,96 e 47,78 ) – ambas acima da média nacional.
Tendo como foco as notas médias dos estudantes agrupados por Organização Acadêmica, notamos que universidades apresentam as maiores médias, ao passo que as piores notas médias foram observadas nos Centros Universitários. 3.1.2 Estatísticas Básicas de Formação Geral A Tabela 3.2 apresenta as estatísticas básicas de Formação Geral por grupo de estudantes.
O total de estudantes inscritos para a prova foi de 5620.
Destes, 77,5 \% do total compareceram, sendo que a presença foi maior entre concluintes ( 88,0 \% ) do que entre ingressantes ( 68,6 \% ).
A média geral da parte da prova de Formação Geral foi de 40,04, sendo que os ingressantes obtiveram média mais baixa ( 39,77 ) que os concluintes, cuja média foi 40,29.
O desvio-padrão geral foi de 20,9, igual para ingressantes e concluintes.
A nota máxima foi 100, obtida por estudantes de ambos os grupos.
Para esta parte, assim como para a nota geral, os coeficientes de assimetria são positivos para todos os grupos.
Como, porém, os valores são bem perto de zero, estes indicam uma ligeira inclinação para a esquerda.
O Gráfico 3.3 permite a comparação entre concluintes e ingressantes relativamente ao desempenho na parte da prova que aborda a Formação Geral.
Nesta parte, se excluídas as provas com nota zero, tal como na prova como um todo, também os concluintes apresentaram desempenho superior, embora seja perceptível um maior equilíbrio.
A mediana dos concluintes ( 35,6 ) foi superior à dos ingressantes ( 34,2 ).
Dos ingressantes, aproximadamente 31,5 \% alcançaram nota superior a 50 ; dentre os concluintes, o percentual dos estudantes nesta faixa de notas é muito semelhante e igual a 31,2 \%.
No Gráfico 3.4 são apresentadas as informações referentes ao desempenho dos ingressantes e dos concluintes na parte de Formação Geral, em diferentes agregações como Região, Categoria Administrativa e Organização Acadêmica.
Considerando -se as notas médias dos estudantes segundo cada Região, observou -se que, tanto no grupo dos ingressantes quanto no grupo dos concluintes, houve diferenças marcantes.
As médias variam, entre os concluintes, de 30,53 na Centro-Oeste a 44,62 na Região Sudeste.
Entre os ingressantes os valores foram de 31,61 na Centro-Oeste a 43,29 no Sudeste.
A maior dispersão de notas dos estudantes de uma dada Região foi encontrado no Centro-Oeste para os concluintes e no Norte para os ingressantes.
Nas Regiões Norte e Centro-Oeste os concluintes apresentaram médias menores do que os ingressantes, a situação inversa aconteceu nas demais Regiões.
Tendo como foco as categorias administrativas, as maiores notas médias foram das instituições estaduais, tanto para concluintes ( 49,63 ) como para ingressantes ( 51,77 ), enquanto as das instituições municipais foram as menores tanto para concluintes ( 31,87 ) e ingressantes ( 31,72 ).
Os estudantes das instituições privadas se apresentam mais homogêneos do que as demais instituições.
Considerando -se o tipo de Organização Acadêmica, Universidades apresentaram os maiores valores para os concluintes e para os ingressantes.
Os Centros Federais de Educação Tecnológica e os Institutos de Educação, Ciência e Tecnologia não estavam representados nesta área. 3.1.3 Estatísticas Básicas do Componente Específico A Tabela 3.3 apresenta as estatísticas básicas referentes ao Componente Específico da prova, tendo sua média geral, de 41,62, muito semelhante ao desempenho da prova de SECRETARIADO EXECUTIVO como um todo cuja média foi 41,19, como apresentado na Tabela 3.1, que por sua vez foi um pouco superior ao desempenho na Formação Geral, que teve média igual a 40,04, conforme mostrado na Tabela 3.2.
Os concluintes apresentaram um desempenho médio ( 43,89 ) superior ao dos ingressantes ( 39,16 ), semelhante ao ocorrido para a nota da Formação Geral, havendo ainda indícios de maior homogeneidade entre as notas médias dos ingressantes do que entre as dos concluintes, como sugeriu o desvio padrão, menor no grupo dos ingressantes do que no dos concluintes.
As notas máximas obtidas por concluintes e ingressantes foram iguais a 93,4 e 88,5, respectivamente.
As distribuições de notas para cada grupo apresentaram simetria, como evidenciado pelos coeficientes de assimetria iguais a zero.
A situação de assimetria é confirmada pelo Gráfico 3.5, que proporciona uma comparação do desempenho entre concluintes e ingressantes em relação ao Componente Específico da prova.
Claramente a distribuição das notas médias dos ingressantes, à direita daquela dos concluintes, evidenciou o percentual de menor concentração de notas médias nas faixas mais baixas para este grupo.
As medianas das distribuições das notas dos ingressantes e dos concluintes foram 38,8 e 44,4, respectivamente, enfatizando o melhor desempenho dos concluintes nesta parte da prova.
Assim como os Gráficos 3.2 e 3.4, o Gráfico 3.6, apresenta uma comparação dos resultados em relação às Regiões, Categoria Administrativa e Organização Acadêmica, agora levando em conta o desempenho dos ingressantes e dos concluintes no Componente Específico da prova.
Tendo como foco as notas médias dos estudantes, em determinada Região, observa -se que, entre os concluintes, as mais elevadas foram encontradas nas Regiões Nordeste ( 44,27 ), Sudeste ( 46,16 ) e Sul ( 44,67 ), valores acima da média nacional para esta parte ( 43,88 ).
Entre os ingressantes, as maiores notas ocorreram nas Regiões : Norte ( 42,53 ) e Sudeste ( 41,49 ) ambas acima da média nacional ( 39,15 ).
Novamente, para a análise segundo a Categoria Administrativa, todas as entidades estão representadas.
As médias mais elevadas, entre os concluintes foram encontradas em instituições de origem federal ( 49,23 ) enquanto entre os ingressantes foram encontradas nas instituições de origem estadual ( 46,50 ).
As notas mais baixas foram encontradas nas instituições municipais, tanto para concluintes ( 39,65 ) como para ingressantes ( 35,95 ).
Quanto à Organização Acadêmica, o comentário é o mesmo já feito para a nota bruta e referente à Formação Geral : só universidades, centros universitários e faculdades são representadas nesta área.
Considerando -se o tipo de Organização Acadêmica, universidades apresentaram os maiores valores, seguidas das faculdades e dos centros universitários.
Assim como os Gráficos 3.2 e 3.4, o Gráfico 3.6, apresenta uma comparação dos resultados em relação às Regiões, Categoria Administrativa e Organização Acadêmica, agora levando em conta o desempenho dos ingressantes e dos concluintes no Componente Específico da prova.
Tendo como foco as notas médias dos estudantes, em determinada Região, observa -se que, entre os concluintes, as mais elevadas foram encontradas nas Regiões Nordeste ( 44,27 ), Sudeste ( 46,16 ) e Sul ( 44,67 ), valores acima da média nacional para esta parte ( 43,88 ).
Entre os ingressantes, as maiores notas ocorreram nas Regiões : Norte ( 42,53 ) e Sudeste ( 41,49 ) ambas acima da média nacional ( 39,15 ).
Novamente, para a análise segundo a Categoria Administrativa, todas as entidades estão representadas.
As médias mais elevadas, entre os concluintes foram encontradas em instituições de origem federal ( 49,23 ) enquanto entre os ingressantes foram encontradas nas instituições de origem estadual ( 46,50 ).
As notas mais baixas foram encontradas nas instituições municipais, tanto para concluintes ( 39,65 ) como para ingressantes ( 35,95 ).
Quanto à Organização Acadêmica, o comentário é o mesmo já feito para a nota bruta e referente à Formação Geral : só universidades, centros universitários e faculdades são representadas nesta área.
Considerando -se o tipo de Organização Acadêmica, universidades apresentaram os maiores valores, seguidas das faculdades e dos centros universitários. 3.2 Análise das Questões Objetivas 3.2.1 Formação Geral A Tabela 3.4 apresenta as estatísticas básicas relativas às oito questões objetivas da Formação Geral dos estudantes.
Adiantamos que a primeira questão foi desconsiderada pelo critério da correlação ponto-bisserial.
Como pode ser observado na Tabela 3.4, os estudantes concluintes obtiveram desempenho superior, com nota média de 50,76, em relação aos estudantes ingressantes cuja nota média foi 50,13.
A variabilidade entre os concluintes foi cerca de 0,1 menor do que entre os ingressantes, como indicaram os desvios-padrão de 22,4, para o grupo de ingressantes e 22,3 para o grupo dos concluintes.
A amplitude total das notas foi a mesma para ambos os grupos : nota mínima de 0 e máxima de 100.
A mediana dos concluintes foi 57,1 e a dos ingressantes, 57,1.
A assimetria, medida pelo coeficiente de assimetria ( skewness ) foi 0,0 para os ingressantes e -0,1 para os concluintes, ou seja, temos uma distribuição muito próxima da simétrica.
A Tabela 3.5 apresenta a classificação das questões objetivas da Formação Geral, segundo o índice de facilidade ( questões em vermelho foram desconsideradas pelo critério ponto-bisserial ).
Como as questões foram comuns a todas as áreas, a tabela apresenta a classificação considerando -se os concluintes de todas as áreas e aqueles do curso de SECRETARIADO EXECUTIVO.
De acordo com os índices obtidos, as questões objetivas da prova foram assim avaliadas para a área de SECRETARIADO EXECUTIVO : das oito questões, duas tiveram entre 60 \%, inclusive, e 85 \% de acertos e foram classificadas como fáceis ( somente uma para o total de concluintes ) ; outras três foram consideradas medianas, alcançando entre 40 \%, inclusive, e 60 \% de acertos ( três para o total de concluintes ) ; duas obtiveram entre 15 \%, inclusive, e 40 \% de acertos, sendo considerada difíceis ( três para o total de concluintes ).
Por fim, uma foi classificada como Muito difícil tento obtido percentual menor do que 15 \% de acertos ( uma para o total de concluintes ), coincidentemente a questão desconsiderada pelo critério ponto-bisserial.
Nenhuma questão foi classificada na categoria Muito fácil.
Em linhas gerais podemos dizer que os concluintes de SECRETARIADO EXECUTIVO se deram melhor na componente de Formação Geral do que o total da população de concluintes de todos os cursos.
A Tabela 3.6 descreve os resultados das análises das questões objetivas relativas à Formação Geral, segundo o poder de discriminação, utilizando -se para tal o índice de discriminação ponto-bisserial.
Como as questões foram comuns a todas as áreas, a tabela apresenta a classificação considerando -se os concluintes de todas as áreas e aqueles do curso de SECRETARIADO EXECUTIVO.
A desconsideração da questão segundo este critério considerou o universo de todos os concluintes.
Uma das oito questões apresentaram índices igual ou acima de 0,40 ( nenhuma para a população de concluintes como um todo ), sendo, portanto, classificadas como muito boa nesta característica, para este grupo de estudantes, duas foram classificadas como boas com índice entre 0,30, inclusive, e 0,40 ( três para os concluintes como um todo ) e quatro foram classificadas como médias, com índice entre 0,20, inclusive, e 0,30 ( quatro para os concluintes como um todo ).
Uma das questões foi classificada como fraca tanto para a população de concluintes como um todo como para a sub-população de concluintes de SECRETARIADO EXECUTIVO e foi desconsiderada pelo critério ponto-bisserial. alcançando entre 40 \%, inclusive, e 60 \% de acertos ( três para o total de concluintes ) ; duas obtiveram entre 15 \%, inclusive, e 40 \% de acertos, sendo considerada difíceis ( três para o total de concluintes ).
Por fim, uma foi classificada como Muito difícil tento obtido percentual menor do que 15 \% de acertos ( uma para o total de concluintes ), coincidentemente a questão desconsiderada pelo critério ponto-bisserial.
Nenhuma questão foi classificada na categoria Muito fácil.
Em linhas gerais podemos dizer que os concluintes de SECRETARIADO EXECUTIVO se deram melhor na componente de Formação Geral do que o total da população de concluintes de todos os cursos.
A Tabela 3.6 descreve os resultados das análises das questões objetivas relativas à Formação Geral, segundo o poder de discriminação, utilizando -se para tal o índice de discriminação ponto-bisserial.
Como as questões foram comuns a todas as áreas, a tabela apresenta a classificação considerando -se os concluintes de todas as áreas e aqueles do curso de SECRETARIADO EXECUTIVO.
A desconsideração da questão segundo este critério considerou o universo de todos os concluintes.
Uma das oito questões apresentaram índices igual ou acima de 0,40 ( nenhuma para a população de concluintes como um todo ), sendo, portanto, classificadas como muito boa nesta característica, para este grupo de estudantes, duas foram classificadas como boas com índice entre 0,30, inclusive, e 0,40 ( três para os concluintes como um todo ) e quatro foram classificadas como médias, com índice entre 0,20, inclusive, e 0,30 ( quatro para os concluintes como um todo ).
Uma das questões foi classificada como fraca tanto para a população de concluintes como um todo como para a sub-população de concluintes de SECRETARIADO EXECUTIVO e foi desconsiderada pelo critério ponto-bisserial.
Na sequência, no Quadro 3.1, estão especificadas, por questão, as habilidades aferidas pelas questões de Formação Geral da prova de SECRETARIADO EXECUTIVO, bem como os resultados da análise relativa aos índices de facilidade e de discriminação de cada questão.
O índice de facilidade a área de SECRETARIADO EXECUTIVO teve uma variação considerável, oscilando entre 0,10 até 0,80 e o de discriminação ficou entre 0,01 e 0,44.
A questão 7, que aferiu as habilidades de “ Analisar e criticar informações ” ; ” Extrair conclusões por indução e/ou dedução ”, ” Estabelecer relações, comparações e contrastes em diferentes situações ”, “ Detectar contradições ”, “ Fazer escolhas valorativas avaliando consequências ”, “ Questionar a realidade ” e “ Argumentar coerentemente ” foi a que apresentou, para a população de concluintes de SECRETARIADO EXECUTIVO, maior poder discriminatório, com índice de 0,44.
A questão mais fácil, com 0,80 de acertos, foi a de número 8, que aferiu também, além das já listadas, a habilidade de “ Ler e interpretar textos ”.
De outra parte, a questão com menor índice de facilidade e de discriminação em Formação Geral foi a questão 1, que foi desconsiderada, e que deveria aferir as habilidades de “ Ler e interpretar textos ”, “ Analisar e criticar informações ” ; ” Extrair conclusões por indução e/ou dedução”,”Estabelecer relações, comparações e contrastes em diferentes situações ”, “ Detectar contradições ”, “ Fazer escolhas valorativas avaliando consequências ”, “ Questionar a realidade ” e “ Argumentar coerentemente ”.
O Gráfico 3.7 mostra, a título de ilustração, o comportamento das respostas à questão 2 de Formação Geral para os estudantes concluintes da área de SECRETARIADO EXECUTIVO como função da nota geral na prova.
Trata -se de uma questão considerada de dificuldade média para os concluintes de SECRETARIADO EXECUTIVO e para o conjunto de concluintes de todas as áreas.
O poder de discriminação foi médio para os estudantes da área e para o conjunto de todos os concluintes.
Neste Gráfico, cada uma das cinco curvas, representou o percentual de respostas em determinada alternativa da questão, em função do escore total dos estudantes concluintes da área de SECRETARIADO EXECUTIVO.
A curva vermelha corresponde à alternativa C, a correta para esta questão.
Assim, observa -se que entre os estudantes com notas mais baixas – desempenho geral mais fraco na prova como um todo –, houve a escolha de alternativas incorretas, especialmente a alternativa E – curva roxa, alternativa D – curva em verde escuro, além da alternativa correta.
A proporção de estudantes que selecionaram a resposta correta aumentou com o desempenho deste grupo de estudantes chegando ao acerto de 100 \% dos estudantes com notas pelo menos 87,5.
Esta análise permitiu verificar como a questão discriminou entre os grupos de desempenho.
Os Gráficos relativos às demais questões de Formação Geral constam do Anexo I. 3.2.2 Componente Específico A Tabela 3.7 apresenta as estatísticas básicas em relação às questões objetivas do Componente Específico.
Como pode ser observado, os estudantes concluintes obtiveram desempenho superior, com uma média de 47,83 pontos, em relação aos estudantes ingressantes com uma média de 43,01 pontos.
Os desvios-padrão indicaram que a variabilidade entre os concluintes foi maior ( 16,2 ) do que entre os ingressantes ( 15,2 ).
Os coeficientes de assimetria para os dois grupos foram praticamente zero indicando uma simetria em volta da média.
A assimetria foi, porém, ligeiramente maior em módulo entre os concluintes.
A análise das questões objetivas do Componente Específico em SECRETARIADO EXECUTIVO, em relação ao grau de facilidade da prova ( ver Tabela 3.8 ), mostrou que, das 27 questões, 12 encontraram -se na categoria difícil ( entre 15 \% e 40 \%, não incluso, de acertos ), 8 foram de nível médio ( entre 40 \% e 60 \%, não incluso, de acertos ), 4 estão na categoria fácil ( entre 60 \% e 85 \%, não incluso, de acertos ) e 1 teve a classificação muito fácil ( a partir de 85 \% de acertos ), 2 questões tiveram um percentual de acertos menor do que 15 \% ( classificação muito difícil ).
Nenhuma questão foi desconsiderada pela Comissão Assessora da Área de SECRETARIADO EXECUTIVO.
Das 27 questões objetivas relativas ao Componente Específico ( ver Tabela 3.9 ), 3 obtiveram o índice de discriminação muito bom,11 obtiveram o índice de discriminação bom e 7 obtiveram nível de discriminação médio.
O índice fraco que eliminaria a questão, foi alcançado por 6 das questões.
Na sequência, no Quadro 3.2 estão especificadas, por questão, as habilidades aferidas pelas questões do Componente Específico da prova, bem como os resultados da análise relativa aos índices de facilidade e de discriminação de cada questão.
O Gráfico 3.8 mostra, a título de ilustração, o comportamento da questão 13 do Componente Específico para os estudantes concluintes da área de SECRETARIADO EXECUTIVO como função da nota geral na prova5.
Trata -se de uma questão considerada de dificuldade média para os concluintes de SECRETARIADO EXECUTIVO.
O poder de discriminação foi bom para os estudantes da área.
Neste Gráfico, cada uma das cinco curvas, representa o percentual de respostas em determinada alternativa da questão, em função da nota total dos estudantes concluintes de SECRETARIADO EXECUTIVO.
A curva roxa correspondeu à alternativa E, a correta para esta questão.
Assim, observou -se que entre os estudantes com escores mais baixos – desempenho geral mais fraco na prova como um todo –, houve a escolha de alternativas incorretas, como a alternativa C – curva em vermelho e a alternativa D – curva em verde escuro, além da alternativa correta.
A proporção de estudantes que selecionaram a resposta correta E, aumentou com o desempenho deste grupo de estudantes chegando ao acerto de 100 \% dos estudantes com notas pelo menos 87,5.
Esta análise permitiu verificar como a questão discriminou entre os dois grupos de desempenho.
Os Gráficos relativos às demais questões do Componente Específico constam do Anexo I. 5 Aqui também, as notas da prova foram grupadas em 8 intervalos de amplitude 12,5, fechados à direita e abertos à esquerda e um ponto com a nota máxima. 3.3 Análise das Questões Discursivas 3.3.1 Formação Geral A Tabela 3.10 apresenta as estatísticas básicas da prova por grupo de estudantes.
A análise dos resultados de desempenho dos estudantes nas questões discursivas que são apresentados na tabela 3.10 e no gráfico 3.9, mostrou que as notas foram mais baixas no conjunto dessas questões que no das questões objetivas.
O total de estudantes inscritos para a prova foi de 5620.
Destes, 77,5 \% compareceram, sendo que a presença foi maior entre concluintes ( 88,0 \% ) do que entre ingressantes ( 68,6 \% ).
A média geral da prova foi de 24,52, sendo que os ingressantes obtiveram média ligeiramente mais baixa, de 24,33, que os concluintes, cuja média foi 24,71.
O desvio-padrão geral foi igual para concluintes e ingressantes, de 33,9.
A nota máxima foi 100,0, atingida por concluintes e ingressantes.
Os coeficientes de assimetria para os dois grupos considerados, o de concluintes e o de ingressantes, bem como para a população como um todo são todos positivos, indicando uma maior concentração à esquerda da distribuição.
Para todos os grupos um pouco mais de 60 \% dos estudantes presentes obtiveram nota zero nas questões discursivas da Formação Geral.
Das 2.688 provas com nota zero, 133 ( 4,9 \% ) continham um protesto, 795 ( 29,6 \% ) estavam em branco e 1760 ( 65,5 \% ) foram respondidas, porém de forma inadequada.
Entre os concluintes a maioria ( 66,9 \% ) dentre aqueles com nota zero foram respondidas, porém de forma inadequada.
Entre os ingressantes 64,0 \% das notas nulas foram respondidas, porém de forma inadequada.
Uma comparação entre o desempenho de ingressantes e concluintes pode ser realizada a partir dos dados apresentados no Gráfico 3.9.
As notas dos concluintes, de um modo geral, foram ligeiramente mais elevadas do que as dos ingressantes, cuja distribuição de notas está à direita dos primeiros.
A ideia de melhor desempenho foi reforçada pela proporção de notas nulas : 61,0 \% entre os concluintes e 62,5 \% entre os ingressantes.
A mediana é igual a 0,0 para ambos os grupos.
Na sequência, discorrer -se-á, especificamente, sobre uma das questões discursivas de Formação Geral, a número 10 já que a número 9 foi desconsiderada.
Em especial, serão analisados os conteúdos englobados na questão e os resultados obtidos pelos estudantes.
Ao final, são apresentados os comentários da banca corretora a respeito do observado nas respostas dos estudantes, suas impressões e conclusões.
Cumpre esclarecer que, tendo em vista que as questões discursivas de Formação Geral foram padronizadas, ou seja, constam de todas as provas, os comentários da banca foram os mesmos para todas as áreas acadêmicas, eis que direcionados a todos os estudantes que participaram do ENADE/2009. 3.3.1.1 Análise da Questão Discursiva 9 de Formação Geral DESCONSIDERADA 3.3.1.2 Análise da Questão Discursiva 10 de Formação Geral A Questão Discursiva nº. 10 buscou avaliar as habilidades : • Ler e interpretar textos ; • Analisar e criticar informações ; • Extrair conclusões por indução e/ou dedução ; • Estabelecer relações, comparações e contrastes em diferentes situações ; • Fazer escolhas valorativas avaliando consequências ; • Questionar a realidade ; • Argumentar coerentemente.
O Quadro 3.3 apresenta as habilidades/competências, temas predominantes e estatísticas básicas da questão 10 de Formação Geral por grupo de estudantes.
Cumpre notar que, como a questão 9 foi desconsiderada, a nota referente às questões discursivas da Formação Geral é idêntica à nota da questão 10.
A média geral da prova foi de 24,52, sendo que os ingressantes obtiveram média ligeiramente mais baixa, de 24,33, que os concluintes, cuja média foi 24,71.
O desvio-padrão geral e igual para concluintes e ingressantes foi de 33,9.
A nota máxima foi 100,0, atingida por concluintes e ingressantes.
Os coeficientes de assimetria para os dois grupos considerados, o de concluintes e o de ingressantes, bem como para a população como um todo foram todos positivos, indicando uma maior concentração à esquerda da distribuição.
Para todos os grupos um pouco mais de 60 \% dos estudantes presentes obtiveram nota zero nesta questão discursiva da Formação Geral.
Das 2.688 provas com nota zero, 133 ( 4,9 \% ) continham um protesto, 795 ( 29,6 \% ) estavam em branco e 1760 ( 65,5 \% ) foram respondidas, porém de forma inadequada.
Entre os concluintes a maioria ( 66,9 \% ) dentre aqueles com nota zero foram respondidas, porém de forma inadequada.
Entre os ingressantes 64,0 \% das notas nulas ( zero ) foram respondidas, porém de forma inadequada.
Uma comparação entre o desempenho de ingressantes e concluintes para esta questão foi exatamente igual à já realizada para a nota discursiva de Formação Geral a partir dos dados apresentados no Gráfico 3.9. 3.3.1.3 Comentários sobre as respostas à Questão Discursiva 10 Apresentação A questão 10 de Formação Geral apresentou um questionamento capaz de analisar a habilidade do estudante em reconhecer o que é um direito social, optar por um direito social vigente que possa diminuir as desigualdades de renda familiar no Brasil e discutir os argumentos em favor do direito social escolhido como suporte à resposta.
Para análise das respostas apresentadas pelos estudantes, houve uma leitura inicial bastante criteriosa.
Foram considerados para efeito de pontuação, os seguintes argumentos : • Acesso à educação pública, gratuita e de qualidade, garantindo ao cidadão melhores condições de vida ; • Assistência à saúde, em seu contexto mais amplo, garantindo os cuidados essenciais de saúde, baseados em técnicas apropriadas ao alcance de todas as pessoas, contando com a participação da população ; • Permanência do estudante na escola, em todos os níveis escolares ( da educação infantil à educação superior, possibilitando a qualificação profissional ), o que garante melhores condições de trabalho e remuneração.
Considerou -se a qualificação e o desempenho docente, desenvolvendo nos estudantes melhores expectativas de aprendizagem e interesse pelo conteúdo da escola.
Os professores desenvolvendo metodologias adequadas e conhecimento pleno quanto a questões de aprendizagem ; • Condições dignas de trabalho, com remuneração que garanta qualidade de vida.
Foram consideradas as condições de trabalho, a preparação profissional para o trabalho e os cursos de aperfeiçoamento, oportunizando momentos de aprendizagem e desenvolvimento, a partir de situações desencadeadoras ; • Ser proprietário do imóvel que reside, reduzindo os gastos com aluguel, garantindo melhores condições de vida, equilíbrio financeiro e construção do bem-estar familiar.
Desempenho dos estudantes na questão É importante salientar que havia respostas que desenvolveram adequadamente o questionamento apresentado, com fundamentação baseada em argumentos coerentes.
Contudo, em um número considerável de respostas, havia apenas um argumento e não dois, conforme fora solicitado.
Muitos estudantes informaram sobre problemas da educação, de um modo geral.
Algumas observações mencionadas pelos corretores durante o processo de correção foram : • há preocupação com uma educação de qualidade, com sugestões de cursos profissionalizantes e métodos preparatórios de aulas. • a educação é revelada como de baixa qualidade, exaltando -se treinamentos e cursos de atualização para a classe docente. • a saúde como essência para a vida.
O SUS é colocado como parte integrante de direito social, garantindo às pessoas fácil acesso. • oportunidade de emprego como forma de ascensão social, direito, cidadania. • os políticos brasileiros como causadores do despreparo educacional.
As respostas analisadas apresentaram norma culta, coesão e coerência razoáveis, que garantiram a compreensão da mensagem desenvolvida pelos estudantes.
Os erros ortográficos foram poucos e não comprometeram a discussão sobre o assunto e sobre os pontos de vista abordados.
Quanto à coerência, houve um desvio quanto ao objetivo da questão, já que alguns estudantes, através dos seus argumentos retrataram reivindicações, reclamações e situações de problemas enfrentados pelos brasileiros, citando golpes políticos.
Houve erros quanto à concordância verbal e nominal.
Houve, também, dificuldade em utilizar os conectores corretamente.
Alguns estudantes, por falta de ideias, apresentaram em suas respostas redundância e prolixidade.
O desempenho foi satisfatório, pois a maioria dos estudantes soube responder o questionamento.
Houve domínio do conteúdo, pois as ideias, os argumentos e a originalidade corresponderam ao esperado. 3.3.2 Componente Específico Nas questões discursivas do Componente Específico – questões de números 38, 39 e 40 –, a diferença de desempenho entre os grupos não foi muito acentuada, inclusive menor do que para as questões discursivas de Formação Geral.
A Tabela 3.11 contém as estatísticas da média destas questões.
Observa -se que no grupo dos concluintes a nota média alcançada foi 21,82, enquanto que no grupo dos ingressantes foi de 17,58, ambas menores do que as obtidas nas questões discursivas de Formação Geral, que foram de 24,3 para os ingressantes e 24,7 para os concluintes ( Tabela 3.10 ).
As médias foram também menores dos que as obtidas para as questões objetivas do Componente Específico ( Tabela 3.7 ) : respectivamente 43,0 e 47,8 para ingressantes e concluintes.
A análise de cada uma destas questões será feita nas subseções seguintes.
Os valores positivos para a medida de assimetria indicaram uma distribuição com cauda para a direita.
Dos 5620 presentes, 33,6 \% obtiveram nota zero.
A tabela 3.12 apresenta um detalhamento da situação destes estudantes.
A maioria ( 55,4 \% ) não fez nenhuma das 3 questões.
Esta proporção foi ligeiramente maior entre os ingressantes ( 56,8 \% ) do que entre os concluintes ( 53,8 \% ).
A combinação de questões em branco e preenchidas correspondeu a 93,9 \% das provas.
Menos de 1 \% no total das provas correspondeu a 3 questões com protesto.
O Gráfico 3.10 apresenta a distribuição das notas nas questões discursivas do Componente Específico por grupo de estudantes.
A moda se situou no primeiro intervalo, aquele com valores entre 0 e 10 incluindo os extremos do intervalo.
A impressão de assimetria informada pelos coeficientes de assimetria positivos ( Tabela 3.11 ) foi reforçada no gráfico pela concentração no primeiro intervalo.
O Gráfico 3.10 apresenta a distribuição das notas nas questões discursivas do Componente Específico por grupo de estudantes.
A moda se situou no primeiro intervalo, aquele com valores entre 0 e 10 incluindo os extremos do intervalo.
A impressão de assimetria informada pelos coeficientes de assimetria positivos ( Tabela 3.11 ) foi reforçada no gráfico pela concentração no primeiro intervalo. 3.3.2.1 Análise da Questão Discursiva 38 do Componente Específico A questão 38, cujas habilidades aferidas encontram -se descritas no Quadro 3.4, que ainda contém outras informações, seguido da Tabela 3.13 e do Gráfico 3.11, foi aquela em que os grupos de estudantes observados obtiveram um desempenho intermediário entre as questões discursivas do Componente Específico.
A nota média do grupo dos ingressantes foi de 17,16, e a dos concluintes 24,89.
O percentual de notas zero obtidas nos dois grupos foi de 56,8 \%, um pouco menor entre os concluintes ( 50,7 \% ) do que entre os ingressantes ( 63,3 \% ).
Mais de metade dos concluintes e dos ingressantes que tiraram zero deixaram a questão em branco ( ver Tabela 3.13 ).
Uma pequena minoria ( 5,3 \% ) utilizou a questão para protestar.
Apesar do baixo desempenho, em ambos os grupos a nota máxima de 100,0 pontos foi alcançada.
O Gráfico 3.11 apresenta a distribuição das notas na questão 38 discursiva do Componente Específico por grupo de estudantes.
A moda se situou no primeiro intervalo, aquele com valores entre 0 e 10 incluindo os extremos do intervalo.
A impressão de assimetria informada pelos coeficientes de assimetria positivos ( Quadro 3.4 ) foi reforçada no gráfico pela concentração no primeiro intervalo. 3.3.2.2 Comentários sobre as respostas à Questão Discursiva 38 Esta questão foi formulada com o intuito de avaliar a habilidade de um futuro secretário executivo em redigir um documento oficial, partindo de uma situação-problema.
Tal meta poderia ser alcançada através do desenvolvimento de um texto coerente, coeso, alicerçado em linguagem culta e formal, porém objetivo, simples e direto.
Não se tratava de uma questão complexa, mas os estudantes apresentaram dificuldades em alcançar o objetivo proposto.
Quanto à abordagem do tema pelos estudantes, pode -se concluir que : A ) A grande dificuldade apresentada pelos estudantes analisados no ENADE/2009 foi a sobre a interpretação de textos e enunciados, não conseguindo deles extrair o que era pedido e/ou esperado.
Por isso, o erro mais comum e gritante desta questão foi a repetição do enunciado.
B ) Muitos estudantes estruturaram o ofício apenas repetindo as informações dadas pela prova com pequenas alterações.
A grande maioria não conseguiu entender que o documento deveria partir do curso de secretariado executivo da Universidade de Pernambuco à Secretaria de Ensino Superior.
C ) Foram poucos os estudantes que conseguiram desenvolver um texto limpo, sem vícios de linguagem ou fuga à norma culta.
Entretanto, dentre esses poucos que atingiram o objetivo proposto, vigora um texto limpo, coerente e eficaz, fugindo, assim, da antiga ideia de que a linguagem rebuscada deva ser empregada em documentos como esse, e atingindo o padrão hoje requerido dos profissionais de secretariado executivo.
D ) Por ser matéria específica do curso, a redação oficial não é ministrada antes que se atinja, pelo menos, a metade do curso.
Mas, entendemos que qualquer graduando, em qualquer área do conhecimento, deixando de lado os aspectos de formatação e diagramação de uma redação oficial, pode e deve ser capaz de elaborar um texto formal solicitando a prorrogação de um prazo estabelecido.
Portanto, as várias provas em branco não devem ser, apenas, justificadas pelo fato de não ter cursado, ainda, tal disciplina.
A vista de tal explicitação em relação à atuação dos estudantes verifica -se que há muito que melhorar em matéria de comunicação e expressão escrita, contudo o mais alarmante diz respeito à capacidade de interpretar textos muito simples que, nesse caso, mostrou -se praticamente nula.
Não deverão existir profissionais nesta área com alcance de uma redação perfeita e que demonstre a compreensão do que se pede ou se declara. 3.3.2.3 Análise da Questão Discursiva 39 do Componente Específico A questão 39, cujas habilidades aferidas encontram -se descritas no Quadro 3.5, que ainda contém outras informações, seguido da Tabela 3.14 e do Gráfico 3.12, foi aquela em que os grupos de estudantes observados obtiveram o desempenho melhor entre as questões discursivas do Componente Específico.
A nota média do grupo dos ingressantes foi de 26,03, e a dos concluintes 29,27.
Isso ocorre, principalmente pelo grande percentual de notas nulas : 42,0 \% entre os ingressantes e 38,1 \% entre os concluintes.
Das provas com notas nulas nesta questão, uma grande proporção foi deixada em branco : 61,8 \% entre os concluintes e 58,3 \% entre os ingressantes ( ver Tabela 3.14 ).
Entre os com nota nula, uma pequena minoria ( 2,9 \% ) utilizou a questão para protestar.
Apesar do baixo desempenho, em ambos os grupos a nota máxima de 100,0 pontos foi alcançada.
Os coeficientes de assimetria positivos indicaram uma concentração à esquerda e espalhamento à direita. documento deveria partir do curso de secretariado executivo da Universidade de Pernambuco à Secretaria de Ensino Superior.
C ) Foram poucos os estudantes que conseguiram desenvolver um texto limpo, sem vícios de linguagem ou fuga à norma culta.
Entretanto, dentre esses poucos que atingiram o objetivo proposto, vigora um texto limpo, coerente e eficaz, fugindo, assim, da antiga ideia de que a linguagem rebuscada deva ser empregada em documentos como esse, e atingindo o padrão hoje requerido dos profissionais de secretariado executivo.
D ) Por ser matéria específica do curso, a redação oficial não é ministrada antes que se atinja, pelo menos, a metade do curso.
Mas, entendemos que qualquer graduando, em qualquer área do conhecimento, deixando de lado os aspectos de formatação e diagramação de uma redação oficial, pode e deve ser capaz de elaborar um texto formal solicitando a prorrogação de um prazo estabelecido.
Portanto, as várias provas em branco não devem ser, apenas, justificadas pelo fato de não ter cursado, ainda, tal disciplina.
A vista de tal explicitação em relação à atuação dos estudantes verifica -se que há muito que melhorar em matéria de comunicação e expressão escrita, contudo o mais alarmante diz respeito à capacidade de interpretar textos muito simples que, nesse caso, mostrou -se praticamente nula.
Não deverão existir profissionais nesta área com alcance de uma redação perfeita e que demonstre a compreensão do que se pede ou se declara. 3.3.2.3 Análise da Questão Discursiva 39 do Componente Específico A questão 39, cujas habilidades aferidas encontram -se descritas no Quadro 3.5, que ainda contém outras informações, seguido da Tabela 3.14 e do Gráfico 3.12, foi aquela em que os grupos de estudantes observados obtiveram o desempenho melhor entre as questões discursivas do Componente Específico.
A nota média do grupo dos ingressantes foi de 26,03, e a dos concluintes 29,27.
Isso ocorre, principalmente pelo grande percentual de notas nulas : 42,0 \% entre os ingressantes e 38,1 \% entre os concluintes.
Das provas com notas nulas nesta questão, uma grande proporção foi deixada em branco : 61,8 \% entre os concluintes e 58,3 \% entre os ingressantes ( ver Tabela 3.14 ).
Entre os com nota nula, uma pequena minoria ( 2,9 \% ) utilizou a questão para protestar.
Apesar do baixo desempenho, em ambos os grupos a nota máxima de 100,0 pontos foi alcançada.
Os coeficientes de assimetria positivos indicaram uma concentração à esquerda e espalhamento à direita.
O Gráfico 3.12 apresenta a distribuição das notas na questão 39 discursiva do Componente Específico por grupo de estudantes.
A moda se situou no primeiro intervalo, aquele com valores entre 0 e 10 incluindo os extremos do intervalo.
A impressão de assimetria informada pelos coeficientes de assimetria positivos ( Quadro 3.5 ) foi reforçada no gráfico pela concentração no primeiro intervalo. 3.3.2.4 Comentários sobre as respostas à Questão Discursiva 39 Esta questão foi considerada como “ Fácil ” e apresentou um facilitador : um texto pronto, sendo necessário apenas reescrevê -lo nos moldes atuais da norma culta.
Das questões discursivas, essa foi a que os estudantes tiveram melhor rendimento.
Um número infinitamente menor de estudantes deixou de responder a essa questão, o que eleva, substancialmente, a precisão da análise proposta pela avaliação.
Considerando os resultados, tivemos casos de repetição do texto tal qual apresentado no enunciado.
Pode -se sintetizar os pontos altos e baixos da correção, tendo em vista o desempenho dos estudantes da seguinte forma : – Alguns estudantes alteraram apenas as palavras de linguagem destacadamente informal, esquecendo -se de posicionar um vocativo ; – Outros estudantes não definiram um sobrenome para o destinatário, citar seu cargo e/ou departamento na empresa. – Muitos estudantes apenas reescreveram o texto ; – Deixaram de prestar atenção a elementos importantes e imprescindíveis à consecução do objetivo proposto.
A partir de tais considerações, torna -se essencial destacar que a capacidade de revisar e reescrever textos parte do pressuposto de que o indivíduo que o faz pode bem interpretar o que leu e, a partir de tal exercício cognitivo, expressar de melhor forma a mensagem previamente escrita.
Considerando tais exposições, conclui -se que os estudantes acabaram ficando muito presos ao texto inicial, ao seu formato e à sua estrutura quando até mesmo isso poderia e deveria ser alterado.
Dessa forma, mostrou -se prejudicada a capacidade de redigir um texto coeso, de \textbf{traçar} essa linha invisível que une as palavras em frases e as frases em parágrafos.
Isto provavelmente fundamentado na dificuldade de interpretação e de domínio de técnicas de reescrita textual, essenciais à profissão.
Muitos erros ortográficos também foram observados nessa questão.
Erros grotescos e infundados quando o público analisado é composto de estudantes ingressantes e concluintes do Ensino Superior.
Concluindo, essa questão deveria ter sido uma das questões mais bem pontuadas pelos estudantes avaliados pelo ENADE no curso de secretariado executivo, não somente por apresentar um grau de dificuldade baixo, mas também por se tratar de atividade rotineira da profissão.
Não há como esperar que um profissional consiga se expressar em outro idioma, quando nem mesmo consegue dominar sua própria língua materna. 3.3.2.5 Análise da Questão Discursiva 40 do Componente Específico A questão 40, cujas habilidades aferidas encontram -se descritas no Quadro 3.6, que ainda contém outras informações, seguido da Tabela 3.15 e do Gráfico 3.13.
A questão foi aquela em que os grupos de estudantes observados obtiveram o pior desempenho entre as questões discursivas do Componente Específico.
A nota média do grupo dos ingressantes foi de 9,60, e a dos concluintes 11,37.
Nesta questão ocorre também o maior percentual de notas nulas : 82,9 \% entre os ingressantes e 79,5 \% entre os concluintes.
Aqui também a maioria das provas com nota zero decorreu de provas deixadas em branco : 72,4 \% entre os ingressantes e 74,6 \% entre os concluintes ( ver Tabela 3.15 ).
Uma pequena minoria ( 4,1 \% ), dentre aqueles com nota nula, utilizou a questão para protestar.
Em ambos os grupos a nota máxima de 100,0 pontos foi alcançada Os coeficientes de assimetria positivos indicaram uma concentração à esquerda e espalhamento à direita.
O Gráfico 3.13 apresenta a distribuição das notas na questão 40 discursiva do Componente Específico por grupo de estudantes.
A moda se situou no primeiro intervalo, aquele com valores entre 0 e 10 incluindo os extremos do intervalo.
A impressão de assimetria informada pelos coeficientes de assimetria positivos ( Quadro 3.6 ) foi reforçada no gráfico pela concentração no primeiro intervalo. 3.3.2.6 Comentários sobre as respostas à Questão Discursiva 40 O objetivo dessa questão era testar a capacidade do estudante de se comunicar através da língua estrangeira, demonstrando, assim, dominar leitura e interpretação instrumentais.
A redação requerida nesse item avaliativo era de formato simples, não exigindo conhecimento e vocabulário avançados.
No entanto, essa foi a questão que apresentou o maior número de abstenções de respostas.
Isso prova que o Inglês ainda continua sendo um tabu entre os estudantes de secretariado executivo.
Uma grande parcela dos estudantes que não responderam à questão alegou não ter atingido ainda o nível requerido de domínio do idioma.
Outra grande parte dos estudantes optou apenas por deixar a questão em branco.
Dentre aqueles que redigiram ( ou tentaram redigir ) o exercício proposto, muitos apenas copiaram o enunciado.
Assim, mais do que nas questões anteriores, esse comportamento se superou em número.
Já dentre os estudantes que realmente responderam à questão pode -se observar grande quantidade de erros gramaticais, tanto de estruturas da língua inglesa quanto em relação à grafia das palavras.
Mais da metade não conseguiu estruturar de forma clara o recado que pretendia passar.
Isso demonstra que sequer houve um entendimento completo do enunciado da questão e muito menos uma redação coesa e coerente.
Devemos ressaltar que para o exercício dessa profissão, é imprescindível o domínio, no mínimo, do inglês instrumental.
Isso inclui até mesmo redações oficiais em outro idioma, tais como cartas, memorandos, ofícios e atas, e não apenas anotações de recados e transmissão de mensagens simples.
Portanto, reitero que, da forma como foi cobrado o conhecimento de inglês no ENADE/2009 era fundamental o alcance de um resultado muito mais satisfatório do que o que foi verificado nas respostas apresentadas.

% matched lemmas: traçar
\end{textsample}
